\chapter{Integration with the Compiler}
\label{sec:compiler}

\section{A pass after type checking}
\label{sec:compiler-pass}

The ownership pass runs where the effect inference runs, after every type is solved and before the
interface is published. Futhark moved its alias checking out of its type checker in 2026 for the same
reason: the old design ``had always been hindered by having to work with incomplete type
information'', whereas a separate pass is ``given a fully well-typed program and simply had to figure
out whether any of the in-place constraints were violated (essentially, whether in-place updates are
semantically observable)'' \citep{henriksen2026rewrite}. The pass never touches unification or the
type store, and its messages carry solved types. It needs three things the checker does not keep
today: the solved type of every list-typed instruction, to decide paths and crossing; a cell-site
identifier per instruction (the instruction index, as the write-set analysis already uses the
instruction to name the indices it tracks); and the source span of every holder, for messages, which
each instruction's main token gives. The write-set analysis already asks the checker for typed paths,
so the hook for the first exists.

\section{Effects}
\label{sec:compiler-effects}

The two effect bits and the ownership classes ride on the same function-type sites, are solved in the
same per-module phase, and are published in the same kind of interface block. They interact in one
place: a function that \emph{suspends} parks its locals in a continuation, which is one-shot (A7), so a
parked local is the same holder it was and liveness is unchanged. A \code{sync} mark says nothing about
ownership. An impure call is a call like any other for holders, and the optimiser's duty to keep
impure calls in place is what A5 extends.

\section{Static dispatch}
\label{sec:compiler-dispatch}

A method call's callee is the term the dispatch table resolves, and its summary is that term's;
evidence a function receives as a parameter is a class instantiated at its caller
(\S\ref{sec:inf-generic}). Derived \code{eq} and \code{compare} are read-only ($\borrowed$). The
inferred \code{where} constraints and the summary block are both inferred facts in the interface, and
both churn the same way; the instrument that measured the first (\S\ref{sec:bg-modules}) can measure
the second.

\section{The write-set analysis}
\label{sec:compiler-writesets}

The direct platform's whole-program write-set analysis (\S\ref{sec:bg-browser}) and the ownership
checker meet only at \code{update}'s entry. There are two reads and one restatement.

\begin{itemize}
\item \emph{The checker computes the entry state itself} (\Obl{P5} in \S\ref{sec:hard-tea}) and reads
from the write-set analysis only the program's shape: which declaration is \code{update}, which is
\code{init}, and the tree of message keys. The write-set analysis does not analyse commands and records
no aliasing between two model paths, so it cannot supply the entry state; the escape scan that the
direct platform would need for in-place records is this same computation, done once for both.
\item \emph{The write-set consumer reads nothing new}, but one word of its soundness statement changes.
Its notion of ``changed'' was ``not the same value, not \code{===}''. Under in-place writes a write
inside \code{model.rows} leaves the list \code{===} to its old self. The write set still \emph{names} the
written path, because it is read off the arm's text rather than off identity, so a consumer that runs
the DOM updates the write set names stays correct; what it may no longer do is conclude ``unchanged''
from \code{===} on a list path, which is A4 and the direct platform's tag guards. The restatement:
\emph{``changed'' means ``may differ in content from what was there before the dispatch''}. The
write-set analysis's own requirement that an in-place list write never touch a list the program still
holds becomes the theorem of Chapter~\ref{sec:soundness}, rather than a requirement on the emitter.
\item \emph{The write-set interpreter is the right engine for the model's paths, not for the general
checker.} It is whole-program, per message, and tracks identity; the ownership checker is per
function, directional over cells, and publishes summaries. They meet at \code{update}'s entry and
nowhere else.
\end{itemize}

\section{The backend}
\label{sec:compiler-backend}

Every accepted demand is in place, so the backend needs no per-site decision: it emits each writer's
call as it does today, and the core library's writers \emph{are} the in-place operations
(Table~\ref{tab:backend}).

\begin{table}[t]
\centering\small
\begin{tabularx}{\linewidth}{@{}>{\raggedright\arraybackslash}p{0.27\linewidth}>{\raggedright\arraybackslash}p{0.3\linewidth}X@{}}
\toprule
Primitive & Today & Under the rule \\
\midrule
\code{set}, \code{update}, \code{swap} & copy up to 256 elements, else trie path copy & \code{a[i] = v} on the base; $O(1)$ \\
\code{push} & copy under 32 elements, else claim the tail & \code{a.push(v)}; amortised $O(1)$ \\
\code{pop} & copy up to 256 elements, else a new header & \code{a.pop()}; $O(1)$ \\
\code{insertAt}, \code{removeAt} & $O(n)$ plain copy & \code{a.splice(i, …)}; $O(n)$ move, as hand-written code does \\
\code{cons} (\code{[ x, …xs ]}) & claim a head slot, short copy, or convert & \code{a.unshift(x)}; $O(n)$ (Chapter~\ref{sec:discussion}) \\
\code{append} & push onto a trie or concatenate & push each element of the second list; $O(|ys|)$ \\
\code{copy} & --- & \code{a.slice()} \\
write through a view $(b, o)$ & --- & the same on $b$ at offset $o + i$; a new header for \code{push}/\code{pop} \\
\bottomrule
\end{tabularx}
\caption{List primitives today and under the rule.}
\label{tab:backend}
\end{table}

\emph{The trie goes.} Its claims, its radix offsets, its head and tail buffers, the cache of its plain
form and its thresholds all exist to make writes to \emph{shared} lists cheap, and under the rule no
shared list is ever written. What stays is the plain form; the view form for a pattern's \code{rest}
(with the scalar-view loop rewrite still removing most views); the array builder used by functions
that build their result in tail position; and the three-fact protocol for readers (\code{length},
\code{Array.isArray}, \code{\$plain()}, the last now needed only for views). Our prototype study measured
the list runtime at 706 bytes compressed without the trie, against 1\,514 with it, and a table
application on the direct platform was found shipping the list-writing half of that runtime for a
list it only reads. All of that is removed. The library's loops written in beni are unchanged.

The backend must still keep the identity promises the library chooses to keep (\S\ref{sec:disc-options},
change~2) and the linearity of the builder. \code{++} on lists lowers to \code{append}, a demand on its
first argument; the generic \code{++} on other appendable types, such as strings, stays a copy
(\S\ref{sec:disc-options}, change~4).

\section{The release optimiser}
\label{sec:compiler-release}

The verdict is a function of the source, and \code{--release} emits the same in-place calls as a
development build, so beni's promise that a release build behaves exactly as a development build holds
with no new exception---\emph{provided the optimiser treats a demand as a barrier} (A5). Today
single-use inlining folds a pure binding into its use across other such bindings, because the language
lets evaluations that make no call commute (\S\ref{sec:bg-release}); a read of a cell and a demand on
it do not commute, so the demand summaries become an input to the optimiser as the impure bit is, and
a read of a cell is never moved across a demand on it. With that, the rest of the optimiser's contract
carries over: an unused binding whose right-hand side is a pure call may be dropped (an unused
\code{List.set} is then not performed, which no reader can tell); two surviving calls are never
reordered; inlining substitutes bodies and never duplicates an argument's evaluation; single-use
inlining may inline a $\mathit{once}$ closure's body into its one call; specialised copies inherit the
general verdict (\S\ref{sec:hard-optimiser}). Renaming and minification touch names, not writes.
Reachability is untouched: \code{copy} ships only when reached.

\section{Records in place}
\label{sec:compiler-records}

Records are outside this thesis's scope, but the machinery applies: a record path is a ``cell'' whose
demand is the in-place record update the lowering would emit, and its holders are the same places.
The language's promise that a record update keeps the other fields as the very same values
(\S\ref{sec:bg-beni}) becomes the only difference: a record written in place keeps its own identity
too, which the direct platform's rule already forbids the renderer to rely on. The hazard is that the
renderer compares records by identity everywhere---a reference check per dynamic hole is what makes
beni's compiled templates fast---so in-place records need a stronger form of A4 than lists do. The
hand-over obligations of \S\ref{sec:hard-tea} are the obligations in-place records would need too.

\section{Incremental builds}
\label{sec:compiler-incremental}

The summary block is in the interface and its hash, and beni's interface firewall fires when it
changes (\S\ref{sec:bg-modules}). Two consequences need measurement. The first is churn: a body edit
that turns a borrowed parameter consumed re-checks every importer, necessarily, since their side
conditions change. The second is the warm rebuild budget---15\,ms for a body edit and 60\,ms for a
signature edit---since a body edit that changes a summary \emph{is} a signature edit under this design,
and signature edits do not yet meet their budget for other reasons.
